\documentclass[a4paper,11pt]{article}
\usepackage[utf8]{inputenc}
\usepackage[T1]{fontenc}
\usepackage[french]{babel}
\usepackage{lmodern}
\usepackage{amsmath,amssymb,amsthm}
\usepackage{mathrsfs}
\usepackage{enumitem}
\usepackage{multicol}

\setlist[enumerate]{label=\textbf{\textcolor{Couleur8}{\arabic*.}}, left=1em}

% Si vous voulez aussi modifier les listes enumerate imbriquées :
\setlist[enumerate,2]{label=\textcolor{Couleur8}{\alph*)}}
\setlist[enumerate,3]{label=\textcolor{Couleur8}{\roman*)}}
\setlist[enumerate,4]{label=\textcolor{Couleur8}{\Alph*)}}
\usepackage{graphicx}
\usepackage{xcolor}
\usepackage{tcolorbox}
\usepackage{geometry}
\usepackage{fancyhdr}
\usepackage{titlesec}
\usepackage{cancel}
\usepackage{ifthen}
\usepackage{tikz}
\usetikzlibrary{arrows.meta}
\usepackage{tkz-tab}
\usepackage{eurosym}

% OPTION PROF/ÉLÈVE
% Décommentez UNE des deux lignes suivantes :
\newboolean{prof}\setboolean{prof}{true}   % Version professeur (avec corrections des devoirs)
%\newboolean{prof}\setboolean{prof}{false}  % Version élève (sans corrections des devoirs)

\definecolor{Couleur1}{HTML}{4A5D7A} %Th
\definecolor{Couleur2}{HTML}{A5B5D6} %Prop
\definecolor{Couleur3}{HTML}{A8C68A} %Méthode
\definecolor{Couleur6}{HTML}{8A9CC1} %Def
\definecolor{Couleur7}{HTML}{E6B459} %Rq Voc
\definecolor{Couleur8}{HTML}{D36740} %Titres
\definecolor{correction}{HTML}{D36740} % Couleur pour les corrections
\definecolor{coopmathscorr}{HTML}{F15929} % Couleur corrections coopmaths

% Configuration des marges
\geometry{
	top=2cm,
	bottom=2cm,
	inner=1.5cm,
	outer=1.5cm,
}

% Police
\renewcommand{\familydefault}{\sfdefault}

% Compteurs
\newcounter{seancenum}
\newcounter{seancenumD}
\setcounter{seancenum}{1}
\setcounter{seancenumD}{1}

% Configuration des en-têtes et pieds de page
\pagestyle{fancy}
\fancyhf{}
\ifthenelse{\boolean{prof}}{
    \fancyhead[L]{\textcolor{Couleur8}{\textbf{Corrections Automatismes - Classe de seconde - Version Professeur}}}
}{
    \fancyhead[L]{\textcolor{Couleur8}{\textbf{Corrections Automatismes - Séance 13 - Classe de seconde}}}
}
\fancyhead[R]{\textcolor{Couleur8}{\textbf{2025-2026}}}
\fancyfoot[L]{\textcolor{Couleur8}{Mme Bertail et Mme Pinto}}
\fancyfoot[R]{\textcolor{Couleur8}{\thepage}}
\renewcommand{\headrulewidth}{1pt}
\renewcommand{\headrule}{\hbox to\headwidth{\color{Couleur8}\leaders\hrule height \headrulewidth\hfill}}

% Environnements personnalisés
\newtcolorbox{seance}[1]{
    colback=Couleur6!10,
    colframe=Couleur1!65,
    fonttitle=\bfseries\large,
    title=Correction Séance \theseancenum #1,
    left=5mm,
    right=5mm,
    top=3mm,
    bottom=3mm,
    arc=0mm,
    after=\stepcounter{seancenum}
}

\newtcolorbox{seanceD}[1]{
    colback=Couleur6!5,
    colframe=Couleur6!45,
    fonttitle=\bfseries\large,
    title=Correction Devoirs - Séance \theseancenumD #1,
    left=5mm,
    right=5mm,
    top=3mm,
    bottom=3mm,
    arc=0mm,
    after=\stepcounter{seancenumD}
}

\begin{document}

\setcounter{seancenum}{13}
\newpage
\begin{seance}{} %13
\begin{enumerate}
\item  $1{,}7\text{ h} =1 \text{ h} + 0{,}7 \text{ h}$, soit $1\text{ h} + \dfrac{7}{10} \text{ h} = 60 \text{ min} + 42\text{ min} = 102 \text{ min}$.\\
          Ainsi, $1{,}7$ heures correspond à ${\color[HTML]{f15929}\boldsymbol{102}}$ {\color[HTML]{f15929}minutes}.


\item Dans une $1$ heure, il y a $3\times 20$ minutes.\\
    Ainsi, en $1$ heure, elle parcourt $20$ fois plus de distance, soit $0{,}9\times 20=18\text{ km}$.\\
    Sa vitesse moyenne est donc ${\color[HTML]{f15929}\boldsymbol{18}}\text{ km/h}$.

\item 
 On remplace $x$ par $x+2$ dans l'expression de $v$ :\\
         $\begin{aligned}
        v(x+2)&=-2(x+2)-10 \\
        &=-2x-4-10\\
        &={\color[HTML]{f15929}\boldsymbol{-2x-14}}
        \end{aligned}$


\item  On se ramène à une équation du type $ax=b$ en isolant les  \og{} $x$ \fg{} dans le membre de gauche et les \og{} non $x$ \fg{} dans le membre de droite.\\
    $\begin{aligned}
    (6x-1)-(4x+4)&=0\\
   6x-1-4x-4&=0\\
 2x-5&=0\\
 2x&=5\\
x&= \dfrac{5}{2}
\end{aligned}$\\
    
   La solution de cette équation est  ${\color[HTML]{f15929}\boldsymbol{\dfrac{5}{2}}}$.


\item  La fonction $f(x) = -3x-6$ est une fonction affine de coefficient directeur $a = -3$ et d'ordonnée à l'origine $b = -6$.\\
    • La fonction s'annule quand $-3x-6 = 0$, soit $x = -2$.\\
    • Comme $a = -3 < 0$, la fonction est décroissante : elle est positive avant la racine et négative après.\\
    Le tableau de signes de la fonction $f$ est donc le suivant :\\
    \begin{tikzpicture}[baseline, scale=0.75]
    \tkzTabInit[lgt=3,deltacl=0.8,espcl=2.1]{ $x$ / 1.5, $f(x)$ / 1.5}{ $-\infty$, $-2$, $+\infty$}
	\tkzTabLine{  , +, z, -}
	\end{tikzpicture}


\item 
 En notant $D$ la dépense totale, Alice a dépensé $\dfrac{1}{9}D$ et Louis, $\dfrac{8}{9}D$.\\
Leur participation équilibrée doit être de $\dfrac{1}{2}D$ chacun.\\
Alice doit donc donner à Louis la somme de $\dfrac{1}{2}D - \dfrac{1}{9}D = \dfrac{7}{18}D$.\\
La fraction de la dépense totale que Alice doit donner à Louis est donc ${\color[HTML]{f15929}\boldsymbol{\dfrac{7}{18}}}$.



\end{enumerate}
\end{seance}

\end{document}